% GENERATED FILE. Edit canonical lessons, then run npm run generate:books.
% canonical-source-hash: fnv1a64:7afa08d905dea129
% canonical-lessons: LA-C492-consolor, LA-C492-consolatio, LA-C492-solacium, LA-C492-lenis, LA-C492-lenio

\chapter{Comfort and gentleness}
\label{ch:la-comfort-and-gentleness}
\hlchaptermodality{492}

\begin{chapteropening}
\textbf{By the end of this chapter:} \emph{I can talk about comforting someone.}

Complete the last lesson of chapter 492: I can talk about comforting someone.
\end{chapteropening}

\section[cōnsōlor, cōnsōlārī]{cōnsōlor, cōnsōlārī — to console, comfort}

\label{lesson:LA-C492-consolor}



\begin{quote}
Before the new one: say the Latin for a groan, then the Latin for to sigh.
\end{quote}

\subsection*{You'll want to know: cōnsōlor, cōnsōlārī}
\textbf{cōnsōlor, cōnsōlārī} — \textquotedblleft{}to console, comfort\textquotedblright{}.

\begin{quote}
Amīca puellam trīstem cōnsōlātur.
\end{quote}

\begin{quote}
\emph{Her friend comforts the sad girl.}
\end{quote}

Deponent: \emph{cōnsōlor, cōnsōlārī, cōnsōlātus sum}, with the accusative. English \emph{console}.

\begin{cousinweb}
\textbf{cōnsōlor} is built on \emph{con-} and \emph{sōlor} \textquotedblleft{}to comfort\textquotedblright{} (Lewis \& Short; EtymDB). English \emph{console} and \emph{consolation} come from it.
\end{cousinweb}

\begin{grammarlens}[title={What you've built}]
One more word for telling, explaining and describing at B1.
\end{grammarlens}

\subsection*{Your turn}
\begin{itemize}
\raggedright
  \item \emph{Say it:} \emph{cōnsōlor, cōnsōlārī}
  \item \emph{Say it:} \emph{cōnsōlor, cōnsōlārī}, once more
  \item \emph{Say it:} a sentence of your own with \emph{cōnsōlor, cōnsōlārī}
  \item \emph{Recall:} say the Latin for a groan, then the Latin for to sigh, then say \emph{cōnsōlor, cōnsōlārī} again
\end{itemize}

\begin{tcolorbox}[breakable,colback=teal!4,colframe=teal!35!black,title={Before you move on}]
What does cōnsōlor, cōnsōlārī mean? (\textquotedblleft{}To console, comfort\textquotedblright{}.) Say it once more.
\end{tcolorbox}
\section[cōnsōlātiō, cōnsōlātiōnis]{cōnsōlātiō, cōnsōlātiōnis — consolation, comfort}

\label{lesson:LA-C492-consolatio}



\begin{quote}
Before the new one: say the Latin for to sigh, then the Latin for to console.
\end{quote}

\subsection*{You'll want to know: cōnsōlātiō, cōnsōlātiōnis}
\textbf{cōnsōlātiō, cōnsōlātiōnis} — \textquotedblleft{}consolation, comfort\textquotedblright{}.

\begin{quote}
Epistula amīcī mihi magna cōnsōlātiō fuit.
\end{quote}

\begin{quote}
\emph{My friend's letter was a great comfort to me.}
\end{quote}

Feminine, third declension. A favourite word of letter-writers.

\begin{cousinweb}
\textbf{cōnsōlātiō} is the action noun of \emph{cōnsōlor} \textquotedblleft{}to console\textquotedblright{}, with \emph{-tiō} (Lewis \& Short). English \emph{consolation} comes from it.
\end{cousinweb}

\begin{grammarlens}[title={What you've built}]
One more word for telling, explaining and describing at B1.
\end{grammarlens}

\subsection*{Your turn}
\begin{itemize}
\raggedright
  \item \emph{Say it:} \emph{cōnsōlātiō, cōnsōlātiōnis}
  \item \emph{Say it:} \emph{cōnsōlātiō, cōnsōlātiōnis}, once more
  \item \emph{Say it:} a sentence of your own with \emph{cōnsōlātiō, cōnsōlātiōnis}
  \item \emph{Recall:} say the Latin for to sigh, then the Latin for to console, then say \emph{cōnsōlātiō, cōnsōlātiōnis} again
\end{itemize}

\begin{tcolorbox}[breakable,colback=teal!4,colframe=teal!35!black,title={Before you move on}]
What does cōnsōlātiō, cōnsōlātiōnis mean? (\textquotedblleft{}Consolation, comfort\textquotedblright{}.) Say it once more.
\end{tcolorbox}
\section[sōlācium, sōlāciī]{sōlācium, sōlāciī — solace, relief}

\label{lesson:LA-C492-solacium}



\begin{quote}
Before the new one: say the Latin for to console, then the Latin for consolation.
\end{quote}

\subsection*{You'll want to know: sōlācium, sōlāciī}
\textbf{sōlācium, sōlāciī} — \textquotedblleft{}solace, relief\textquotedblright{}.

\begin{quote}
Librī mihi in senectūte sōlācium sunt.
\end{quote}

\begin{quote}
\emph{Books are a solace to me in old age.}
\end{quote}

Neuter, second declension. English \emph{solace} comes from it through Old French.

\begin{cousinweb}
\textbf{sōlācium} is built on \emph{sōlor} \textquotedblleft{}to comfort, soothe\textquotedblright{} (Lewis \& Short; EtymDB). English \emph{solace} comes from it through Old French.
\end{cousinweb}

\begin{grammarlens}[title={What you've built}]
One more word for telling, explaining and describing at B1.
\end{grammarlens}

\subsection*{Your turn}
\begin{itemize}
\raggedright
  \item \emph{Say it:} \emph{sōlācium, sōlāciī}
  \item \emph{Say it:} \emph{sōlācium, sōlāciī}, once more
  \item \emph{Say it:} a sentence of your own with \emph{sōlācium, sōlāciī}
  \item \emph{Recall:} say the Latin for to console, then the Latin for consolation, then say \emph{sōlācium, sōlāciī} again
\end{itemize}

\begin{tcolorbox}[breakable,colback=teal!4,colframe=teal!35!black,title={Before you move on}]
What does sōlācium, sōlāciī mean? (\textquotedblleft{}Solace, relief\textquotedblright{}.) Say it once more.
\end{tcolorbox}
\section[lēnis, lēne]{lēnis, lēne — gentle, mild, soft}

\label{lesson:LA-C492-lenis}



\begin{quote}
Before the new one: say the Latin for consolation, then the Latin for solace.
\end{quote}

\subsection*{You'll want to know: lēnis, lēne}
\textbf{lēnis, lēne} — \textquotedblleft{}gentle, mild, soft\textquotedblright{}.

\begin{quote}
Vōx mātris lēnis et dulcis est.
\end{quote}

\begin{quote}
\emph{Mother's voice is gentle and sweet.}
\end{quote}

Two endings: \emph{lēnis} for masculine and feminine, \emph{lēne} for neuter. Of voices, winds and people.

\begin{cousinweb}
\textbf{lēnis} is a native Latin adjective; Lewis \& Short compares \emph{lentus} \textquotedblleft{}pliant, slow\textquotedblright{}, and EtymDB lists a Proto-Indo-European parent, but the link is not settled.
\end{cousinweb}

\begin{grammarlens}[title={What you've built}]
One more word for telling, explaining and describing at B1.
\end{grammarlens}

\subsection*{Your turn}
\begin{itemize}
\raggedright
  \item \emph{Say it:} \emph{lēnis, lēne}
  \item \emph{Say it:} \emph{lēnis, lēne}, once more
  \item \emph{Say it:} a sentence of your own with \emph{lēnis, lēne}
  \item \emph{Recall:} say the Latin for consolation, then the Latin for solace, then say \emph{lēnis, lēne} again
\end{itemize}

\begin{tcolorbox}[breakable,colback=teal!4,colframe=teal!35!black,title={Before you move on}]
What does lēnis, lēne mean? (\textquotedblleft{}Gentle, mild, soft\textquotedblright{}.) Say it once more.
\end{tcolorbox}
\section[lēniō, lēnīre]{lēniō, lēnīre — to soothe, soften}

\label{lesson:LA-C492-lenio}



\begin{quote}
Before the new one: say the Latin for solace, then the Latin for gentle.
\end{quote}

\subsection*{You'll want to know: lēniō, lēnīre}
\textbf{lēniō, lēnīre} — \textquotedblleft{}to soothe, soften\textquotedblright{}.

\begin{quote}
Medicus dolōrem aegrī medicāmentō lēnit.
\end{quote}

\begin{quote}
\emph{The doctor soothes the patient's pain with a medicine.}
\end{quote}

\emph{lēniō, lēnīre, lēnīvī, lēnītum}, from \emph{lēnis}. English \emph{lenient} comes from its present participle.

\begin{cousinweb}
\textbf{lēniō} is built on \emph{lēnis} \textquotedblleft{}gentle, soft\textquotedblright{} (Lewis \& Short). English \emph{lenient} comes from its present participle \emph{lēniēns}.
\end{cousinweb}

\begin{grammarlens}[title={What you've built}]
One more word for telling, explaining and describing at B1.
\end{grammarlens}

\subsection*{Your turn}
\begin{itemize}
\raggedright
  \item \emph{Say it:} \emph{lēniō, lēnīre}
  \item \emph{Say it:} \emph{lēniō, lēnīre}, once more
  \item \emph{Say it:} a sentence of your own with \emph{lēniō, lēnīre}
  \item \emph{Recall:} say the Latin for solace, then the Latin for gentle, then say \emph{lēniō, lēnīre} again
\end{itemize}

\begin{tcolorbox}[breakable,colback=teal!4,colframe=teal!35!black,title={Before you move on}]
What does lēniō, lēnīre mean? (\textquotedblleft{}To soothe, soften\textquotedblright{}.) Say it once more.
\end{tcolorbox}
